% =============================================================================
%  1. Panorama: qué es, qué hace hoy, qué quedó afuera
% =============================================================================

\begin{frame}{Para qué es este documento}
  \begin{block}{Objetivo}
    Tener en un solo lugar \destaca{todo el trabajo}: qué hace el robot, cómo
    está armado, qué hace cada archivo y cada parámetro, y \destaca{de dónde
    salió cada número}. Es un resumen para mí, no para proyectar: tiene más
    texto que una presentación normal.
  \end{block}

  \vspace{0.3em}
  Cómo está organizado:
  \begin{enumerate}
    \item \textbf{Panorama} --- qué hace hoy y qué quedó afuera.
    \item \textbf{Hardware} y \textbf{arquitectura} --- las piezas y cómo se hablan.
    \item \textbf{Archivo por archivo} --- el código del Pi y el firmware del ESP32.
    \item \textbf{Comportamiento} --- la máquina de estados y la ley de control.
    \item \textbf{Parámetro por parámetro} --- qué hace cada uno y por qué vale lo que vale.
    \item \textbf{Historia} --- las 15 jornadas, lo que se midió y lo que falló.
    \item \textbf{Operación y pendientes} --- cómo se usa y qué falta.
  \end{enumerate}

  \vspace{0.2em}
  {\footnotesize Fuentes: \cod{PLAN.md} (el plan y la bitácora), \cod{HARDWARE.md},
  \cod{RUNBOOK.md}, los comentarios de \cod{src/config.py} y los registros de
  \cod{logs/}. Si algo de acá no coincide con esos archivos, mandan ellos.}
\end{frame}

\begin{frame}{Índice}
  \setbeamertemplate{section in toc}[sections numbered]
  \begin{columns}[T,onlytextwidth]
    \begin{column}{0.48\textwidth}
      {\footnotesize\tableofcontents[sections={1-8},hideallsubsections]}
    \end{column}
    \begin{column}{0.48\textwidth}
      {\footnotesize\tableofcontents[sections={9-15},hideallsubsections]}
    \end{column}
  \end{columns}
\end{frame}

% =============================================================================
\section{Panorama}
% =============================================================================

\begin{frame}{El objetivo, en una frase}
  \begin{alertblock}{}
    \centering\large
    Un robot que \textbf{recorre una sala}, \textbf{reconoce dos objetos con la
    cámara} y \textbf{se comporta distinto con cada uno}: se acerca a uno y huye
    del otro.
  \end{alertblock}

  \vspace{0.6em}
  \begin{columns}[T]
    \begin{column}{0.48\textwidth}
      \begin{block}{Objetivo: el oso (\cod{teddy bear})}
        Un oso \textbf{impreso en una hoja A4} ($\sim$20 cm). El robot lo busca,
        se le acerca corrigiendo el rumbo con la cámara y frena a $\sim$30 cm.
      \end{block}
    \end{column}
    \begin{column}{0.48\textwidth}
      \begin{block}{Amenaza: la mochila (\cod{backpack})}
        Una mochila real. El robot retrocede, gira media vuelta y se aleja.
        Tiene \textbf{prioridad incondicional} sobre todo lo demás.
      \end{block}
    \end{column}
  \end{columns}

  \vspace{0.6em}
  \small
  Todo corre \textbf{a bordo}: una Raspberry Pi 5 con una red neuronal
  (YOLOv8n) decide, y un ESP32 mueve los motores. No hay PC externa, ni
  teleoperación, ni nube. La notebook sólo sirve para lanzar la corrida por SSH
  y mirar lo que pasa.
\end{frame}

\begin{frame}{Qué hace el robot hoy (medido en el piso el 1 y el 2 de octubre)}
  \relax{\footnotesize
  \begin{tabular}{L{2.4cm}L{8.3cm}L{3.0cm}}
    \toprule
    \textbf{Qué ve} & \textbf{Qué hace} & \textbf{Estado} \\
    \midrule
    Nada &
      Mira quieto $\sim$0,6 s, gira un paso de $\sim$18\grad{} y repite. Una
      vuelta entera tarda $\sim$17 s; después avanza $\sim$18 cm para cambiar el
      punto de vista y vuelve a empezar. &
      \cod{BUSCAR} \\[2pt]
    El oso &
      Se queda quieto hasta confirmarlo (3 frames, $\sim$0,3 s), se acerca
      centrándolo, frena a $\sim$30 cm y da un pitido largo. Si lo pierde, lo
      busca hacia el lado donde lo vio por última vez. &
      \cod{APROXIMAR} \hacia{} \cod{ENCONTRADO} \\[2pt]
    La mochila &
      Se queda quieto hasta confirmarla y, con la alarma sonando, retrocede
      $\sim$30 cm girando, sigue girando hasta darle la espalda y avanza
      $\sim$65 cm. Después vuelve a buscar. &
      \cod{HUIR} \\[2pt]
    Algo a menos de 30 cm adelante &
      Cuando va a avanzar sin el oso a la vista (huyendo o buscando), no
      avanza: gira en el lugar hasta tener el camino libre. &
      \cod{HUIR} o \cod{BUSCAR} \\[2pt]
    Oso y mochila &
      Huye: la amenaza gana siempre, incluso estando detenido frente al oso. &
      \cod{HUIR} \\
    \bottomrule
  \end{tabular}}

  \vspace{0.6em}
  \begin{columns}[T,onlytextwidth]
    \begin{column}{0.485\textwidth}
      \begin{block}{Lo que se probó}
        {\footnotesize Dieciséis corridas con motores en dos días: diez el 1
        de octubre y seis el 2, ya con el ultrasónico y el buzzer. Hay dos
        videos vistos desde el robot en \cod{media/}.}
      \end{block}
    \end{column}
    \begin{column}{0.485\textwidth}
      \begin{alertblock}{Lo que no se midió todavía}
        {\footnotesize Las 10 corridas por escenario que piden las métricas, y
        la autonomía del cargador. Los porcentajes de éxito no existen aún.}
      \end{alertblock}
    \end{column}
  \end{columns}
\end{frame}

\begin{frame}{Los números del sistema, de un vistazo}
  \relax{\footnotesize
  \begin{columns}[T,onlytextwidth]
    \begin{column}{0.49\textwidth}
      \begin{tabular}{L{3.9cm}L{2.9cm}}
        \toprule
        \textbf{Percepción} & \\
        \midrule
        Modelo & YOLOv8n (COCO) \\
        Tamaño de inferencia & 320 px \\
        FPS del lazo completo & $\sim$10,6 ($\sim$9,9 grabando) \\
        Inferencia por frame & $\sim$88 ms \\
        Edad del frame al decidir & $\sim$150 ms \\
        Alcance con el oso A4 & $\sim$1,3 m \\
        Campo visual horizontal & $\sim$52\grad \\
        Confirmar / perder & 3 / 5 frames \\
        \bottomrule
      \end{tabular}

      \vspace{0.5em}
      \begin{tabular}{L{3.9cm}L{2.9cm}}
        \toprule
        \textbf{Seguridad} & \\
        \midrule
        Vigencia del comando (Pi) & 0,40 s \\
        Watchdog (ESP32) & 500 ms \\
        Ultrasónico: esquivar / llegar & 30 / 20 cm \\
        Corte de emergencia & conector de 12 V \\
        \bottomrule
      \end{tabular}
    \end{column}
    \begin{column}{0.49\textwidth}
      \begin{tabular}{L{3.9cm}L{2.9cm}}
        \toprule
        \textbf{Movimiento} & \\
        \midrule
        Rango de velocidad & 0,091 a 0,380 m/s \\
        Velocidad máx. usada & 0,25 m/s \\
        Giro buscando & pasos de $\sim$18\grad \\
        Giro en régimen & 124 a 165\grad/s \\
        Distancia de frenado & $\sim$30 cm del oso \\
        Asimetría de los motores & 2,11 veces \\
        Trocha (regla / efectiva) & 21,5 / 18 cm \\
        Rueda & 60 mm de diámetro \\
        \bottomrule
      \end{tabular}

      \vspace{0.5em}
      \begin{tabular}{L{3.9cm}L{2.9cm}}
        \toprule
        \textbf{Verificación} & \\
        \midrule
        Verificaciones sin hardware & 79 \\
        Corridas registradas & CSV + JSON c/u \\
        \bottomrule
      \end{tabular}
    \end{column}
  \end{columns}}
\end{frame}

\begin{frame}{Qué entra en esta versión y qué quedó afuera}
  \begin{columns}[T,onlytextwidth]
    \begin{column}{0.485\textwidth}
      \begin{exampleblock}{Adentro, y funcionando}
        \begin{itemize}
          \item Detección con YOLOv8n a bordo, a $\sim$11 FPS.
          \item Máquina de estados con arbitraje: buscar, aproximar,
                encontrado, huir.
          \item Control proporcional sobre la imagen.
          \item ESP32 con protocolo serie, compensación de zona muerta por
                rueda y watchdog.
          \item Alimentación portátil con dos fuentes separadas.
          \item \textbf{Ultrasónico} al frente: no avanza contra lo que
                tenga a menos de 30 cm. \textbf{Buzzer}: alarma y llegada.
          \item Registro de cada corrida y análisis automático.
          \item 79 verificaciones del comportamiento sin hardware.
        \end{itemize}
      \end{exampleblock}
    \end{column}
    \begin{column}{0.485\textwidth}
      \begin{alertblock}{Afuera: segunda versión}
        \begin{itemize}
          \item \textbf{Sensores atrás y a los costados}: el robot
                \ojo{retrocede y gira a ciegas}, y una pared muy de costado no
                la ve.
          \item \textbf{Encoders}: sin odometría ni lazo de velocidad.
          \item \textbf{Frenar a 10 cm}: a esa distancia el oso no entra en el
                cuadro; haría falta un avance final a ciegas.
          \item \textbf{Cámara CSI} (Arducam): nunca fue vista por el Pi 5; se
                reemplazó por una webcam USB.
        \end{itemize}
      \end{alertblock}
    \end{column}
  \end{columns}

  \vspace{0.5em}
  {\footnotesize El ultrasónico y el buzzer habían quedado afuera el
  2026-10-01, por falta de tiempo, y entraron al día siguiente (sección ``El 2
  de octubre'').}
\end{frame}

\begin{frame}{Cómo se cubren los requisitos del módulo}
  \relax{\footnotesize
  \begin{tabular}{L{4.2cm}L{9.6cm}}
    \toprule
    \textbf{Requisito} & \textbf{Cómo se cubre} \\
    \midrule
    Aplicación de IA &
      Red convolucional de detección de objetos (YOLOv8n, preentrenada en COCO)
      haciendo inferencia en tiempo real sobre hardware embebido. \\[2pt]
    Vehículo autónomo no tripulado &
      Robot terrestre de tracción diferencial, con la decisión 100 \% a bordo. \\[2pt]
    Percepción \hacia{} decisión \hacia{} actuación &
      La cadena completa cerrada en el robot físico: cámara, red, máquina de
      estados, ley de control, enlace serie, motores. \\[2pt]
    Toma de decisiones &
      Máquina de estados con arbitraje explícito de conflictos: huir tiene
      prioridad sobre aproximarse. \\[2pt]
    Fusión de sensores &
      Cámara y ultrasónico. La cámara dice qué hacer y hacia dónde; el
      ultrasónico impide avanzar contra un obstáculo y puede dar por alcanzado
      el objetivo. Entró el 2026-10-02. \\
    \bottomrule
  \end{tabular}}

  \vspace{0.6em}
  \begin{block}{Lo que el trabajo demuestra}
    {\small Integrar \textbf{percepción con IA} en un \textbf{lazo de control
    físico} que corre en \textbf{tiempo real} sobre \textbf{hardware real}. La
    mitad de lo interesante salió de la última parte: motores con zonas
    muertas y asimetrías, imágenes que se rompen, baterías que no sostienen
    los 5 V.}
  \end{block}
\end{frame}

\begin{frame}{Las seis decisiones de diseño (1 de 2)}
  \relax{\small
  \begin{block}{D1 --- Un solo microcontrolador de bajo nivel: ESP32}
    Se descartó usar ESP32 y PSoC5LP a la vez: dos firmwares, dos protocolos y
    dos cadenas de herramientas duplican la superficie de fallas sin sumar
    nada. Se programa con PlatformIO, en estilo C.
  \end{block}
  \begin{block}{D2 --- Modelo preentrenado en COCO, sin entrenamiento propio}
    El ciclo captura \hacia{} etiquetado \hacia{} entrenamiento agregaba
    riesgo de cronograma sin ser el aporte del trabajo, que es \textbf{integrar
    la inferencia en un lazo de control físico}. Se compensa con más
    profundidad en control, arbitraje y medición.
  \end{block}
  \begin{block}{D3 --- Clases por nombre: \cod{teddy bear} y \cod{backpack}}
    Objetos medianos, compactos y de buen contraste. Las clases se declaran
    \textbf{por nombre} y se resuelven contra \cod{model.names}: si cambia el
    modelo, falla al arrancar en vez de detectar otra cosa en silencio.
    \ojo{En el piso apareció un matiz:} la mochila real es \cod{suitcase} para
    el modelo (ver la sección de la jornada final).
  \end{block}}
\end{frame}

\begin{frame}{Las seis decisiones de diseño (2 de 2)}
  \relax{\small
  \begin{block}{D4 --- Huir tiene prioridad sobre aproximar}
    Ante las dos clases a la vez, gana la conducta de seguridad. Es la regla
    clásica de la robótica de comportamientos: los de preservación subsumen a
    los de objetivo.
  \end{block}
  \begin{block}{D5 --- Encoders como extensión, no como núcleo}
    El sistema tiene que andar en lazo abierto (PWM directo, calibrado por
    rueda). Un PID sobre una planta descompensada arranca peor, así que la
    calibración había que hacerla igual. Los encoders quedaron sin cablear.
  \end{block}
  \begin{block}{D6 --- Confirmación temporal de las detecciones}
    Una detección aislada no cambia el estado: hacen falta \textbf{3 frames
    seguidos} para confirmar una clase y \textbf{5 sin verla} para darla por
    perdida. Evita que un falso positivo dispare una maniobra.
  \end{block}}

  \vspace{0.2em}
  {\footnotesize Y dos decisiones de arquitectura que salieron al implementar:
  la compensación de zona muerta vive \textbf{en el ESP32} y no en el Pi, y
  las capas que \emph{deciden} no tocan hardware (ver ``Arquitectura'').}
\end{frame}
